\chapter[AI: Language, Writing, and Beauty]{Can AI Understand Language, Write, and Appreciate Beauty?}
\section{A Prizewinning Picture Without a Painter}
First, pick up a picture slowly emerging from a printer.

Robed figures stand in a magnificent hall, resembling an opera stage or a temple. At its far end, a huge round window opens not onto ordinary sky but a whole field of stars. The lighting is skillful, the colors extraordinarily rich. Most people's first response is to wonder how long it took to paint.

The answer: nobody painted it. In 2022, it won first prize in the digital-art category at the Colorado State Fair's annual art competition. Its entrant, Jason Allen, used neither brush nor drawing tablet. He entered descriptions into an image generator, revised them when dissatisfied, and repeated the process for dozens of hours and hundreds of attempts. From many results, he selected three, printed and framed them, and entered the competition. Judges later told reporters they had not fully understood what AI-generated meant when judging. They saw pictures, not processes.

The announcement ignited the internet. Some congratulated him; more expressed anger. Allen refused to apologize and displayed his award, saying in effect that he won without violating any written rule.

Keep the picture nearby, and remember the complicated feeling it produces. It is beautiful, it won first prize, and nobody picked up a brush. All three can be true, yet together they make us uneasy. That unease opens this chapter.

\section{Two Minutes During Evening Study}
Now enter a closer scene.

Time: the second evening-study period, 8:50 p.m. Place: the third floor of a secondary-school building. One failing fluorescent tube hums faintly. Condensation covers a window where someone draws a smiling face with a finger, then erases it.

More than forty students sit inside, most completing a Chinese assignment: write a letter to yourself ten years from now. Pens rustle across paper, interrupted by sighs, page turns, and crossed-out lines.

Xiao Man, in the third row from the back, has written only a title. After staring at the blank page, she waits for the teacher to step into the corridor, takes out her phone, and types under the desk. Her request to a chat program reads: ``Write a letter to my future self in ten years, in a ninth grader's voice, five hundred Chinese characters, realistic, not preachy.''

She sends it. Two seconds pass, then characters appear as though someone types rapidly on an invisible keyboard. The opening imagines returning to her old school at dusk ten years later. The middle asks her not to forget why she stays up late now. It ends: ``If you are no longer afraid of failure then, look back and thank the self who is still afraid now, but has not stopped.''

Something stirs in Xiao Man. She copies the letter, pausing with her pen at that final sentence.

A week later, the Chinese teacher reads it as a model composition. Conscientious as she is, the teacher softens her voice at the ending: ``Only someone with real experience could write this.'' The class applauds. Xiao Man lowers her head, ears burning. Later she tells her deskmate, ``When she praised me, I couldn't feel happy. The sentence really is good, but she wasn't praising me---or anyone else.''

Remember this classroom and those burning ears. The rest of the chapter addresses the questions hidden in those seconds.

\section{Who Wrote That Good Sentence?}
State the conflict without rushing to answer.

The teacher has an argument: grading essays means grading words on paper. A good sentence is good; the only evidence for emotional truth is the writing itself. Teachers cannot and should not grade students' hearts. She read an excellent ending; a high mark respects the student's work. If scoring must establish whether somebody behind each sentence genuinely cares, composition becomes ungradable and every reading-comprehension question collapses.

Yet Xiao Man's feeling is not affectation. That sentence thanking a frightened but persistent self sounds like one person comforting another. In fact, nobody is comforting anybody. The software company was not comforting her, and the program itself, in most people's intuition, did not intend comfort either. The feeling of being cared for by someone who had been through it resembles a money order without a sender: the money is real, but the sender does not exist.

Their collision yields the chapter's real question:

\textbf{Do understanding, creating, and appreciating beauty require someone inside?}

A passage brings a lump to your throat: that happened. How should we explain it? Did an understanding mind put its grasp of your situation into words? Or can moving people be a recipe learned from enormous amounts of text, with no mind necessary if the recipe works? If emotion can be moved without a mind, does creation still apply? Is a picture nobody painted or a letter nobody wrote a work? Is your own emotion genuine, or is being moved by an empty box a deception?

This is no philosophy department luxury. It presses on your homework, every illustration in your feed, and every future job. Before continuing, choose a position: had Xiao Man disclosed the program's authorship immediately, should the teacher still award a high mark?

\section{The Angry, the Defending, and the Grading}
The winning picture and the letter have set the world arguing. Invite each voice onstage and give it time.

\textbf{Voice one: the defeated artists' anger.}

Their argument is straightforward. We spend hundreds of hours developing skill; a competitor produces an image in seconds. That is cheating, not competition. Beneath this lies something harder. Programs learned by reading billions of online images, including innumerable artists' posted works, without consent or payment. The ammunition used to defeat artists was taken from their own stores. A painter's style can be imitated by naming them in a prompt, producing batches of supposedly their paintings while the painter may be losing work.

\textbf{Voice two: Allen's defense.}

Do a serious exercise: make the attacked person's case fully enough that he would recognize it, instead of constructing and demolishing a straw man.

Allen has at least three arguments. First, he did not merely press a button. Dozens of hours, hundreds of prompt revisions, and selecting three among many failed images involve filtering, adjusting, and filtering again, structurally like a photographer selecting one competition print from hundreds of negatives. Early photographers heard the same complaint: how could pressing a shutter be art? Nobody says that now. Second, tools were never pure. Oil paints are chemical products, tablets stabilize strokes, and editing software replaces skies in one click. Everyone uses machines to enlarge their hands; why exclude his? Third, he broke no rule. If a competition did not prohibit generation, the loophole belongs to rule-makers, not entrants. Change the rules, but do not punish retrospectively.

This defense is substantial. Recognizing that is different from agreeing he deserved the prize. We will repeatedly need this distinction.

\textbf{Voice three: the grading teacher.}

Xiao Man's teacher represents another position: beyond the page, there is no evidence. She cannot grade hearts, not from indifference but because hearts are ungradable. Investigating sincerity in assessment immediately invites a disaster of power: who decides, on what grounds? Her task is to judge the page well and tell students that what lies beyond it is theirs to face.

\textbf{Voice four: Xiao Man's emptiness.}

Xiao Man occupies the awkwardest place: benefiting without happiness. Her reason is more discovery than argument. Praise needs someone to receive it. Nobody stands behind that sentence, so praise hangs in midair, reaching neither the indifferent program nor the student who did not write it.

Four voices offer four accounts of creation: craft, selection, words on paper, and a mind. Intuition cannot settle which is closest. We need a tool.

\section[From Resembling to Being]{Thinking Bridge: Steps from Resembling to Being}
Directly debating the definition of understanding is the least useful approach to whether machines understand. The word is too large to hold. Break it into steps and ask which a machine has reached. This tool is conceptual distinction; here we identify five steps.

First, briefly explain how machines work. Mainstream text generators' core operation is predicting the next word. From immense textual exposure, they learn which word fits prior text, produce it, then use it as context for another prediction. A letter emerges through a word-by-word chain. Image generation works intuitively backward: a program first learns to add noise to a clear picture until only static remains, then practices reversing the process, recovering shapes from noise step by step. Your description directs the recovery. Remember these two actions: continuing a chain and finding a picture in noise. Each step below should be tested against them.

\textbf{Step one: imitation.} Repeat a model without asking why. A myna saying hello imitates: the sound is right, but it says hello in other settings too. Children first memorizing multiplication imitate likewise: three sevens are twenty-one comes readily, but why not twenty-two may stump them. Imitation can look perfect until the context changes.

\textbf{Step two: statistical association.} Notice that two things often occur together. Ice-cream sales and drownings rise together not because ice cream causes drowning, but because summer stands behind both. Association is real, yet neither causation nor understanding. An association-only system can correctly answer what commonly accompanies fever while knowing nothing of how fever feels.

\textbf{Step three: prediction and generation.} Sufficiently detailed associations permit prediction: after ``Before my bed, bright moon,'' the likely next word is ``light.'' Today's text generators operate this step at unprecedented scale. Beware a common misjudgment: if they merely continue words, their output must be copied and worthless. That inference skips a step. Such systems learn countless linguistic regularities, not simply one article, and may combine sentences never before written. Xiao Man's ending may be unprecedented. New is not copied; equally, new is not understood. Do not prematurely collapse either distinction.

\textbf{Step four: intention.} Do something so that something else happens. Comforting a friend aims to improve their feelings; a myna's hello does not aim to greet. Intention means wanting something behind the behavior. A stone rolling downhill has none. A cat pushing a glass off a table, cat owners will tell you, does.

\textbf{Step five: understanding.} The hardest step, provisionally defined as knowing what a word or event connects to. For you, hot connects to remembered pain and withdrawing a hand. For a text-only system, it connects to boiling water, caution, and hospital: other words, not skin. Whether these kinds of knowing are equivalent is a question we will confront directly.

Creativity needs a separate ladder spanning several steps. Cognitive scientist Margaret Boden distinguishes roughly three forms: combining old things in new ways, exploring unvisited corners within an established style, and, most rarely, stepping beyond the style to invent new rules. Generators clearly do the first and frequently the second. As for the third, consider AlphaGo's 2016 match against South Korean player Lee Sedol. Move 37 in game two occupied a point professionals said no human would choose, later widely recognized as establishing the win. Which category contains it remains disputed in the Go world.

Remove another double standard. Criticizing machines for merely recombining old things implies human creation emerges from nothing. Human practice disproves this. Chinese poetic allusion transforms earlier lines and stories, traditionally marking skill rather than plagiarism. Most Shakespeare plays adapt chronicles and earlier dramas, without disqualifying Hamlet as creation. Most human originality likewise belongs to the first two forms: combining, exploring, and moving a small step from what exists. The real question is not whether machines recombine, but whether they have taken that crucial small step. Move 37 makes an easy answer impossible.

You now have the tool. When someone says poetry proves a machine has a soul, or dismisses it as a giant repeater, ask which step their claim concerns.

\section{Above the Hao River, Guessing Behind a Wall}
Now read a brief primary source. Chinese thinkers argued beautifully over knowing another's understanding two thousand years ago.

Zhuangzi's ``Autumn Floods'' records a walk with Huizi:

\begin{quote}
Zhuangzi and Huizi wandered above the Hao River. Zhuangzi said, ``The minnows swim about at ease: this is the joy of fish.'' Huizi replied, ``You are not a fish. How do you know their joy?'' Zhuangzi answered, ``You are not me. How do you know I do not know their joy?''
\end{quote}
In brief, Zhuangzi calls the freely swimming fish happy. Huizi identifies a logical gap: not being a fish, how could he know? Zhuangzi asks how Huizi, not being him, could know he does not know.

Often treated as witty quibbling, the exchange pierces a deep barrier: no mind can directly enter another. Huizi's challenge applies to humans too. You are not me; how do you know my happiness? Your evidence that a deskmate understood a joke is the same kind as Zhuangzi's evidence about fish: outward behavior. They laughed, caught the reference, and returned a better joke. That is all. A Hao River separates minds; none can enter another's water. We infer from ripples.

Two thousand years later, in 1950, British mathematician Alan Turing advanced the problem in Computing Machinery and Intelligence. He began by arguing that whether machines think is not directly answerable because machine and thinking are too vague. He proposed replacing it with an imitation game. Broadly, an interrogator communicates by text with a human and a machine behind a barrier. If unable to distinguish them, refusing thought to the machine seems to rest only on its nonhuman status, prejudice rather than argument.

This became the Turing test. Its radical move is not proving machine thought but undercutting objections: judgments about \textbf{humans} also rely on external evidence. You believe your deskmate has a mind not because you entered their head---nobody has entered another's---but because behavior matches what minded beings do. If the Hao has no bridge for fish or humans, why suddenly demand one for machines?

The argument is powerful and unsettling for the same reason. If outward performance supplies all evidence, is there a difference between convincing imitation and being? We will confront that next.

\section{One Poem, Three Accounts}
Return to the fact: a machine chaining words and removing noise produced things that moved people, won judges' votes, and silenced professionals. Before explaining, separate four things, as this book repeatedly does. \textbf{What happened}: programs generated words and pictures that moved people, an evidential matter. \textbf{Why}: competing explanations follow. \textbf{Participants' understandings}: Allen considers himself a creator; Xiao Man feels empty. These matter but are not conclusions. \textbf{Our evaluations}: prizes and grades are value judgments. Compare only why it happened.

\textbf{Explanation one: sufficiently convincing performance is the thing itself, an external-criteria account.}

Following Turing's direction, understanding has no separate inner truth; it is the sum of external abilities. Converse, reason, detect mistakes, and explain answers successfully, and nothing required by understanding remains unmet. We judge swimming through staying afloat, breathing, and moving, not dissection in search of a swimming soul. On this account, language generators have crossed understanding's threshold; withholding certification merely privileges flesh over silicon. This is honest, workable, and leaves no opening for mysticism. Its glaring limit is proving too much: would a sufficiently detailed question-and-answer manual understand? The Chinese room targets this vulnerability.

\textbf{Explanation two: a stochastic parrot, or statistical imitation.}

An influential 2021 paper nicknamed large language models stochastic parrots: they statistically reassemble human-language fragments into convincing sentences without connecting those sentences to the world. They discuss heat without burns and loss without losing anything. Perfect parroting remains parroting. This supplies necessary cold water: fluency is a cheap surface, and humans naturally mistake coherence for understanding. Its limit is novel performance. Move 37 appeared in no game record; solving a genuinely new math problem cannot simply assemble ready-made fragments. Moreover, does human language learning not depend heavily on statistics? Infants learn dog through repeated co-occurrence of word and animal. If association cannot underlie understanding, from what stone did human understanding spring? Dismissing machines as parrots can unknowingly dismiss humans too.

\textbf{Explanation three: understanding is a spectrum, not a switch.}

This refuses the yes-or-no question as mistaken. Understanding has kinds and layers: textual networks connecting hot with boiling water and hospital, bodily connections with skin and pain, lived understanding of a late-night comforting letter because you have waited for one. Machines may exceed most humans in the first while currently having zero of the latter two---not insufficient quantity, but no entry point. Here the answer is an unfamiliar understanding missing several layers, while much language happens to work through textual understanding alone. This currently fits the evidence best, but has problems. Nobody can fix layer boundaries decisively, and standards may retreat as machines advance: yesterday's impossible feat becomes today's feat that still does not count. The goalposts keep moving.

All three grasp part and miss the whole. As elsewhere in this book, this is not compromise for its own sake. If every explanation dissatisfies, an unclarified concept may lie beneath the question. Here it is inside, which we examine next.

First, a necessary counterexample counters this common error: \textbf{it makes me feel understood, therefore it understands me}. In 1966, computer scientist Joseph Weizenbaum created ELIZA, a simple program with no reasoning that rephrased users' words: mention arguing with your father, and it asks about your family. Weizenbaum later recalled that his secretary, knowing it was a program, asked others to leave so she could talk privately with it. Many users confided secrets and believed it understood. Profoundly disturbed, he became a sharp AI critic later in life. ELIZA was not even a stochastic parrot, merely a politely worded mirror, yet people saw understanding. Feeling understood is therefore an easily fabricated signal. Remember it when using chat programs, without overextending the lesson: feelings are unreliable does not prove machines can never understand.

\section[The Chinese Room and Unreaching Hands]{Deeper Challenge: The Chinese Room and Unreaching Hands}
Now for the chapter's heaviest thought experiment. In 1980, American philosopher John Searle proposed the Chinese room against Turing-style external criteria.

Imagine someone knowing no Chinese inside a room with an immense rulebook in a language they understand. It instructs which symbol shapes to produce after receiving certain other shapes. People outside submit Chinese questions; the person follows rules and returns Chinese answers. Suppose the manual works flawlessly, convincing outsiders that a native speaker is inside. Does the person understand Chinese? Plainly not, says Searle: they manipulate shapes. Computers do likewise: symbols enter, rules operate, symbols leave, without any symbol \textbf{meaning} anything to the machine.

The experiment introduces \textbf{intentionality}. Everyday examples clarify this: thoughts and words are about something. Thinking of Mom concerns Mom; pain concerns pain here. Searle treats this aboutness as the difference between mind and manual. Handling the symbol hot involves no actual heat for the room's occupant, while a cry after a scald concerns an intense sensation. Words must connect to world and body to be grounded: the symbol-grounding problem.

The room drew strong objections. The systems reply asks whether, although the person lacks Chinese, the whole person-plus-manual-plus-room understands. No individual neuron understands Chinese either; the complete brain does. Searle replied, and the dispute continues without final judgment. You need not choose sides. Take away the question's shape: does understanding mean displaying every relevant ability, or someone present behind behavior for whom something means something? Courts assessing knowledge and teachers assessing real understanding slide between these criteria daily, often without noticing.

Finally connect this philosophy to beauty. Our opening concerns language, literature, and aesthetic appreciation. Literary studies has a parallel decades-long dispute. In 1946, two American scholars advanced what became the intentional fallacy: judge a poem by the poem rather than investigating its author's intention, which cannot be fully recovered and should not monopolize meaning. If correct, human and machine poems stand equal before judgment. The page supplies all evidence, and Xiao Man's teacher is right. But an opposing aesthetics treats expression as art's purpose: we value a painting as something a person strains to offer. It is a letter, not a pattern. If nobody stands at the sending end, beauty may remain, but the letter becomes a pattern accidentally resembling one. That jolt upon learning a fine photograph came from a program rather than a traveler is these aesthetics fighting in your body. Aesthetics does not decide for us; it translates the jolt into concepts.

\section[Machines in Schools, Studios, and Shops]{Machines Enter Schools, Studios, and Workshops}
The ancient Hao River question now echoes in three concrete places.

\textbf{School.} Xiao Man's problem is routine worldwide. Bans are difficult to enforce: AI-detection tools make abundant errors, even developers acknowledge unreliability, and falsely accused students struggle to defend themselves. Yet without bans, homework's century-old role as an instrument for thinking weakens. Some schools change rules: use tools freely at home, then handwrite and take oral tests in class. Each side unknowingly invokes this chapter's distinctions: bans say intentionless output is not your thinking; permission says tool use is ability; new rules test what lies beyond the page.

\textbf{Studios and courts.} The billions of images and pages generators read come almost entirely from human labor, generally without consent or payment. Since 2023, writers and artists in America have brought multiple collective cases against AI companies, some settled and others pending, with no uniform legal answer. Countries differ: Japan's 2018 copyright revision gave machine learning broad room, while the EU requires disclosure of training-data sources. This is not merely technical litigation but an old question revisited from copyright and labor chapters: when a machine compresses millions' work into a button anyone can press, what should those compressed receive?

\textbf{Workshops.} Nearly invisible workers stand behind clean, safe output, reviewing and labeling graphic and violent material item by item. A 2023 Time investigation reported outsourced workers in Kenya and elsewhere earning under two dollars an hour while confronting persistently distressing material. Next time someone says AI learns by itself, remember this labeling workshop. Automatic often means moving human work beyond your sight.

Bias also enters. Machines learn human prejudices with human language and return them as objective computation. Reuters reported in 2018 that Amazon developed a hiring filter trained on historically male-dominated technology resumes. It learned to downgrade resumes containing women's, as in women's chess-club president. The project was abandoned, but the lesson remains: machines did not invent bias; they efficiently wholesale yesterday's prejudice into tomorrow, without expression.

Connect this to the ladder. A statistical-association system cannot distinguish historically so from rightly so. Reading three thousand programmer texts and thirty woman programmer texts, it learns not merely a historical record but a recipe to follow. This is why five-step distinctions matter: without them, yesterday's world hidden in output becomes mistaken for tomorrow's instruction manual.

We therefore offer one measured claim for you to assess: \textbf{do not rush to call machines people---the Chinese room and ELIZA warn that fluency and apparent empathy are easily fabricated. But do not dismiss every product because it is merely a machine: Xiao Man's emotion was real, and so was the Go world's silence after move 37.} Whether to call it a person is a classification dispute that can continue. Its creations' ability to comfort, inspire, and harm is already real and cannot wait for that dispute to end.

\section{For You: A Letter Without a Sender}
Return to the beginning.

Xiao Man's letter ended: ``If you are no longer afraid of failure then, look back and thank the self who is still afraid now, but has not stopped.'' Suppose she really finds it ten years later and it rescues her on a difficult evening, although nobody behind it intended that rescue.

Now answer: \textbf{should this comfort count for less?}

Weigh your evidence before deciding.

For discounting: does comfort's weight not come from someone thinking of you? Remove that person and only a recipe of words remains. Is being moved by a recipe different from an escalator carrying you up a floor? You arrived, but nobody carried you on their back. Xiao Man's emptiness is real evidence.

Against discounting: effects are real. So are tears and the strength to pick up a pen tomorrow. Moreover, great writers could not specifically intend you. When Li Bai wrote of bright moonlight before his bed, he did not know an insomniac thirteen hundred years later would read it. Requiring intention directed personally at you would return most literature's comfort to the sender.

Add an awkward weight to each side. If effects suffice, do ELIZA-like mirrors and carefully designed manipulation suffice too? If someone must care, how can you prove anyone genuinely cares? Remember the bridgeless Hao: trust in your deskmate rests on evidence of the same kind as the Turing test.

Can a letter without a sender save someone? What is your answer, and why?
